\chapter{反方与风险：翻倍情景如何失败}

\section{最强反方观点}

最强反方不是“这些公司没有增长”，而是市场可能正确地拒绝为周期利润支付高倍数。雅化与盛新的锂盐利润受商品价格、库存和套保影响；恩捷的2027利润依赖价格、利用率与客户验证；江波龙的存储利润可能处于涨价与低价库存共振的峰值；中国船舶的高价订单兑现虽然更长，但仍属于周期交付。

若分母回落，低PE会迅速变成高PE；若分母兑现但倍数不升，股价也无法在五个月内翻倍。年末目标必须同时承受盈利风险与估值风险。

\begin{exhibitbox}[Exhibit 11：概率—影响风险矩阵]
\small
\begin{tabularx}{\exhibitboxwidth}{L{3cm} C{1.4cm} C{1.4cm} L{2.8cm} X}
\toprule
风险 & 概率 & 影响 & 暴露标的 & 触发与动作 \\
\midrule
锂价反转 & 中高 & 高 & 雅化、盛新、天华 & 立即下修EPS与倍数，删除纯价格驱动Bull \\
利润不转现金 & 高 & 高 & 雅化、江波龙、盛新、天华 & H2 OCF仍负则降级或维持零权 \\
隔膜恢复停滞 & 中 & 高 & 恩捷 & 价格/利用率/良率未达标则下修2027E \\
存储峰值回落 & 高 & 高 & 江波龙 & 毛利与库存转差则拒绝高PE \\
船舶交付或毛利延迟 & 中 & 中高 & 中国船舶 & 下修2027EPS与整合信用 \\
风险偏好压缩 & 高 & 高 & 全部 & 延长兑现期，不沿用年末牛市目标 \\
目标期限未披露 & 高 & 中 & 雅化、恩捷 & 仅雅化通过资格后限权10\%；恩捷零权 \\
\bottomrule
\end{tabularx}
\end{exhibitbox}

\section{现金流是共同薄弱环节}

江波龙2026Q1经营现金流-28.75亿元、盛新锂能-6.67亿元、雅化集团-4.31亿元、天华新能-1.83亿元。雅化与江波龙FY2025经营现金流也为负。盈利上行而现金下行，可能来自库存、预付款、应收、采购保障或扩产；在周期反转时，这些项目会放大风险。

恩捷与中国船舶的现金表现相对更好。恩捷FY2025与Q1均为正，但扩产资本开支仍需覆盖；中国船舶Q1经营现金流81.52亿元，但船舶建造现金流受里程碑回款影响，单季数据不能永久外推。

\section{证据风险与卖方选择偏差}

本报告归档的14份原始券商报告全部为正面评级，没有中性或卖出报告。这不是市场真实评级分布，而是公开可获得报告的选择偏差。雅化与恩捷各只有一份完整点目标；恩捷的关键参数尤其集中于单一来源。

因此，报告不以“券商都看多”作为结论，而是把原始报告拆成可审计的预测、方法、目标与风险，再与公司公告和现金流交叉。目标价是一个假设包，而不是外部真理。

\section{期限风险}

东吴证券两份报告未披露目标价期限；其公司评级定义不能替代具体目标期限披露。42元和103元是否针对12月31日无法确认。将目标价直接贴上“年末”标签，会夸大确定性。

\section{结论失败的明确信号}

若雅化H2经营现金流不转正或利润主要来自库存，唯一条件核心资格取消；若恩捷无法提供客户涨价、单平利润与利用率证据，103元继续保持零权敏感性；若中国船舶不能保持毛利与现金，质量备选下调；若江波龙H2利润与现金继续背离，任何翻倍区间删除。

\begin{riskbox}[风险纪律]
本报告不设置机械止损或交易指令。研究上的“止损”是证据失效：当盈利桥、现金桥或估值期限不再成立，应更新结论，而不是等待股价替研究者证明原叙事正确。
\end{riskbox}

\section{组合压力测试}

压力测试不是预测，而是检验结论能否承受同时冲击。若锂价从当前中枢下跌25\%，雅化与盛新的盈利分母和估值倍数可能同步下修；若隔膜涨价延后两个季度，恩捷的2027E EPS与市场前置定价都需下调；若全市场高弹性板块PE再压缩30\%，即使盈利不变，所有年末翻倍情景都将显著后移。

\begin{exhibitbox}[Exhibit 12：三类压力与结论变化]
\small
\begin{tabularx}{\exhibitboxwidth}{L{3cm} L{3.6cm} X X}
\toprule
压力情景 & 首要受影响变量 & 结论变化 & 恢复条件 \\
\midrule
锂价中枢下跌25\% & 雅化/盛新/天华吨利、库存与PE & 雅化取消条件核心；锂链Bull删除 & 自供降本与现金流抵消价格冲击 \\
隔膜恢复延后两季度 & 恩捷单平利润、利用率、2027E EPS & 103元零权敏感性继续不进入正式模型 & 客户涨价与新线良率形成公司证据 \\
高弹性PE压缩30\% & 全部标的市场倍数 & 年末翻倍概率接近个位数 & 盈利上修足以抵消倍数收缩 \\
存储毛利回落15个百分点 & 江波龙利润与库存价值 & 观察区间下修，禁止使用峰值EPS & 现金转正、库存周转与产品结构改善 \\
\bottomrule
\end{tabularx}
\end{exhibitbox}

\clearpage
